\documentclass[12pt]{article}
\usepackage[brazil]{babel}
\usepackage[utf8]{inputenc}
\usepackage[T1]{fontenc}
\usepackage{amsmath, amssymb}
\usepackage{geometry}
\usepackage{graphicx}
\usepackage{xcolor}
\usepackage{enumitem}

% Definições de Cores
\definecolor{meuazul}{RGB}{0, 79, 151}

\geometry{margin=2.5cm}

\begin{document}

\begin{center}
    \textcolor{meuazul}{\Large \textbf{Lista de Exercícios 02: Segurança da Informação (Parte 2)}}
\end{center}

\vspace{0.8cm}

\begin{enumerate}[label=\textbf{Questão \arabic*:}, leftmargin=2.2cm, labelsep=0.5cm]

    \item \textbf{Definição e Objetivos da Criptografia} \\
    A criptografia é uma ferramenta fundamental na engenharia de segurança. Explique, com suas palavras, o que é o processo de criptografia e quais são seus principais objetivos no contexto da proteção de dados.

    \item \textbf{Componentes do Processo Criptográfico} \\
    Diferencie os seguintes termos fundamentais:
    \begin{enumerate}[label=\alph*)]
        \item Texto Simples (\textit{Plaintext});
        \item Algoritmo Criptográfico;
        \item Texto Cifrado (\textit{Ciphertext});
        \item Chave (\textit{Key}).
    \end{enumerate}

    \item \textbf{Criptografia Simétrica} \\
    A criptografia simétrica, ou de chave privada, é amplamente utilizada devido à sua eficiência.
    \begin{enumerate}[label=\alph*)]
        \item Qual é a principal característica do uso de chaves nesse método?
        \item Cite uma vantagem e a principal desvantagem (desafio) da criptografia simétrica.
    \end{enumerate}

    \item \textbf{Criptografia Assimétrica} \\
    Diferente da simétrica, a criptografia assimétrica utiliza um par de chaves (pública e privada).
    \begin{enumerate}[label=\alph*)]
        \item Explique a função de cada uma das chaves no processo de envio de uma mensagem segura.
        \item Por que a criptografia assimétrica costuma ser combinada com a simétrica em aplicações reais (como no protocolo HTTPS)?
    \end{enumerate}

\end{enumerate}

\end{document}
